\Needspace{15\baselineskip}\fancyhead[L]{\small JUNO. ELEMENTS VII.}
\centerline{JUNO. ELEMENTS VII.}

\begin{center}
\rule{0.5\linewidth}{0.4pt}\\[2pt]
Göttingische gelehrte Anzeigen. Issue 14. Pp. 129, 130. 23 January 1808.\\[2pt]
\rule{0.5\linewidth}{0.4pt}
\end{center}

\par\indent The only observations of \emph{Juno} from last year that have become known are those Inspector \textsc{Bessel} made in Lilienthal in March, April, and May with the ring micrometer. The planet was so faint that it probably could not be observed anywhere with fixed instruments. Moreover, while it was visible, most astronomers were so occupied with the newly discovered \emph{Vesta} that Juno may have been somewhat neglected. Our Professor \textsc{Gauss} therefore did not hesitate to use those observations to improve the latest elements of Juno, especially since Mr. \textsc{Bessel} observed with
\SourceLeaf{298}an excellent telescope and is remarkably skilled in using the ring micrometer. The May observations were first used to determine the opposition, with the following result:

\[
\begin{array}{l}
1807.\ \text{May}\ 17.\quad 12^{\text{u}}\,27^{\text{m}}\,41^{\text{s}}\ \text{M.\,T. in Göttingen}\\[2pt]
\text{True longitude}\ .\ .\ .\ .\ .\ 236^\circ\ 4'\ 57''6.\\[2pt]
\text{True geocentric latitude}\quad 16\;56\;0.1\ \text{north}
\end{array}
\]\par\indent This opposition was then combined with those of 1804 and 1806 and with other positions from 1804, 1805, and 1806 to calculate the following new elements (VII.), which represent all the observations used as closely as possible. [They appear on page 294 of this reprint as part of the letter to \textsc{Bode} of 24 January 1808.]\par\indent The ephemeris for the next apparition of Juno, based on these new elements, has already been completed, and astronomers will find it in Baron \textsc{von Zach}’s Monatliche Correspondenz. This calculation was all the more necessary because the ephemeris Professor \textsc{Bode} constructed in the Yearbook for 1810 from the VI. elements gives right ascensions more than $\frac{3}{4}$ of a degree too large around the time of the next opposition, and Juno will still be so faint this time that finding and observing it will be very difficult.
